Reinforcement learning (RL) can give a congestion controller a learned control law in place of a hard-coded one, and a decade of work has produced many designs for the Internet, wireless, satellite, IoT and data-center networks. We report a systematic survey of this literature following PRISMA~2020. From \nIdentifiedTotal{} records retrieved from OpenAlex, arXiv, USENIX proceedings and forward and backward citation searching, we include \nStudies{} primary studies (\nReportsIncluded{} reports) published from 2016 to October 2026. We code each study's setting, control point, state, action, reward, algorithm, training environment, highest evaluation level, generalization testing, fairness evaluation and code availability, and we release every decision with its evidence.
Designs have shifted from clean-slate neural policies to hybrids in which learning tunes, selects or bounds a classic algorithm: \nPctHybridPost\% of designs since 2023, against \nPctHybridPre\% before. The evidence base is weak and is not improving. Among studies whose evaluation environment is known, \nPctSimOnlyOfKnown\% are evaluated only in simulation and \nPctRealOfKnown\% reach a real network, and the real-network share has fallen since 2021. Only \nPctGen\% of design studies test outside their training distribution, and \nPctCodeDesign\% release code. A reproducible case study shows a PPO controller that achieves \nCsNarrowRefUtil\% utilization in distribution but drops to \nCsNarrowJitterTenUtil\% under 10\,ms of RTT jitter. Broader training helps only within the range it covers, and it made training less reliable. The deployments that exist either keep a classic algorithm or safeguard in the loop or learn from the operator's own logs, are run by operators who own the endpoints, and optimize a narrow service metric. We map nine deployment barriers to the evidence and mitigations in the corpus, review recent work that uses large language models to design and run controllers, and propose an evaluation checklist and a research agenda.
